% BHU PYQ -- B.A./B.Sc. (Hons) Semester IV Examination 2025-2026 (NEP)
% Paper: MATMJ-41 : Calculus of Several Variables
% Exam: B.A./B.Sc. (Hons) Semester IV Examination, 2025-2026 (NEP)

\documentclass[11pt,a4paper]{article}
\usepackage[a4paper,top=2.5cm,bottom=2.5cm,left=2.8cm,right=2.8cm,headheight=14pt]{geometry}
\usepackage{amsmath,amssymb}
\usepackage[T1]{fontenc}\usepackage[utf8]{inputenc}
\usepackage{lmodern,microtype,fancyhdr,enumitem,hyperref,lastpage,parskip}

\hypersetup{colorlinks=true,urlcolor=black,linkcolor=black,
  pdftitle={MATMJ-41 Calculus of Several Variables -- BHU B.Sc. Sem IV 2025-26 (NEP)}}
\pagestyle{fancy}\fancyhf{}
\renewcommand{\headrulewidth}{0.4pt}\renewcommand{\footrulewidth}{0.4pt}
\fancyhead[L]{\small\textit{Banaras Hindu University}}
\fancyhead[R]{\small\textit{MATMJ-41: Calculus of Several Variables (2025-26 NEP)}}
\fancyfoot[C]{\small Page \thepage\ of \pageref{LastPage}}

\newlist{parts}{enumerate}{2}
\setlist[parts,1]{label=(\alph*),leftmargin=*,itemsep=5pt,topsep=3pt}
\setlist[parts,2]{label=(\roman*),leftmargin=*,itemsep=3pt,topsep=2pt}
\newcommand{\pts}[1]{\hfill\textit{[#1]}}

\begin{document}
\begin{center}
  {\large\bfseries BANARAS HINDU UNIVERSITY}\\[2pt]
  {\large\bfseries B.A./B.Sc. (Hons) Semester IV Examination, 2025-26 (NEP)}\\[4pt]
  {\normalsize Subject: Mathematics \hfill Paper No. MAT-MJ-41 : Calculus of Several Variables}\\[4pt]
  \rule{0.8\linewidth}{0.5pt}\\[6pt]
  {\normalsize Time: 2½ Hours \hfill Full Marks: 70}\\[2pt]
  {\small REF NO. Conf/Even/N/2025-26/64}
\end{center}
\rule{\linewidth}{0.4pt}
\smallskip
\noindent\textit{Instructions: Answer \textbf{five} questions including Question 1 which is compulsory. The figures in the right-hand margin indicate marks.}
\bigskip

\noindent\textbf{Question 1.} \textit{(Compulsory)}
\begin{parts}
  \item Give the $\varepsilon - \delta$ definition of simultaneous limit of a real valued function in $\mathbb{R}^3$.\pts{2}
  \item Define Euclidean norm in $\mathbb{R}^n$.\pts{2}
  \item Give the definition of Gamma function.\pts{2}
  \item Show that the mixed partial derivatives of a real valued function in $\mathbb{R}^2$ need not be equal, with the help of an example.\pts{3}
  \item Give the definition of homogeneous function with an example.\pts{3}
  \item Define stationary point of a real valued function in $\mathbb{R}^2$.\pts{2}
  \item Define interior and boundary points of a set in the Euclidean space $\mathbb{R}^n$.\pts{2}
  \item Define triple integral in Cartesian coordinates.\pts{2}
  \item Give the statement of Liouville's Extension of Dirichlet's Theorem.\pts{2}
  \item Define Jacobian of an implicit function.\pts{2}
\end{parts}
\medskip\rule{\linewidth}{0.2pt}

\medskip\noindent\textbf{Question 2.}
\begin{parts}
  \item If $x, y \in \mathbb{R}^n$ and $a \in \mathbb{R}$, then prove that---
  \begin{align*}
    \|x + y\| &\le \|x\| + \|y\| \\
    \|ax\| &= |a| \cdot \|x\|
  \end{align*}\pts{6}
  \item State and prove mean value theorem for a function $f : D \subseteq \mathbb{R}^2 \to \mathbb{R}$.\pts{6}
\end{parts}
\medskip\rule{\linewidth}{0.2pt}

\medskip\noindent\textbf{Question 3.}
\begin{parts}
  \item Investigate for simultaneous limit and continuity of the function
  \[
    f(x,y) = \begin{cases}
      xy \frac{x^2 - y^2}{x^2 + y^2} & ; (x,y) \neq (0,0), \\
      1 & ; (x,y) = (0,0),
    \end{cases}
  \]
  at $(0,0)$.\pts{6}
  \item Prove that a function which is differentiable at a point admits of first-order partial derivatives at the point but the converse is not true.\pts{6}
\end{parts}
\medskip\rule{\linewidth}{0.2pt}

\medskip\noindent\textbf{Question 4.}
\begin{parts}
  \item Find the equation of the tangent plane to the surface $2x^2 y + y^2 z = 3$ at the point $(-1, 1, 1)$.\pts{6}
  \item State and prove Young's theorem.\pts{6}
\end{parts}
\medskip\rule{\linewidth}{0.2pt}

\medskip\noindent\textbf{Question 5.}
\begin{parts}
  \item Expand $f(y, x) = e^{xy}$ in power of $(x-1)$ and $(y-1)$ by Taylor's theorem as far as terms of third degree.\pts{6}
  \item Establish necessary and sufficient conditions for a function $f : D \subseteq \mathbb{R}^2 \to \mathbb{R}$ to have maximum or minimum value at a point $(a,b) \in D$.\pts{6}
\end{parts}
\medskip\rule{\linewidth}{0.2pt}

\medskip\noindent\textbf{Question 6.}
\begin{parts}
  \item Find the volume under the plane $x + y + z = 6$ and above the triangle in the $xy$-plane bounded by $x = y$, $y = 0$ and $x = 3$.\pts{6}
  \item State and prove Euler's theorem for the function of $n$ variables.\pts{6}
\end{parts}
\medskip\rule{\linewidth}{0.2pt}

\medskip\noindent\textbf{Question 7.}
\begin{parts}
  \item By using Dirichlet's integral, find the volume of the solid bounded by $x = 0$, $y = 0$, $z = 0$ and the surface
  \[
    \left(\frac{x}{a}\right)^{1/2} + \left(\frac{y}{b}\right)^{1/2} + \left(\frac{z}{c}\right)^{1/2} = 1
  \]
  where $0 < c < b < a < \infty$.\pts{6}
  \item State and prove Legendre's duplication formula.\pts{6}
\end{parts}

\bigskip
\begin{center}
  $\star \star \star$
\end{center}
\end{document}
